% reality/coordinate_systems-DE.tex

\chapter{Koordinatensysteme und World-Einbettung}
\label{chap:coordinate-systems}

\index{Koordinatensystem}
\index{Bezugssystem!World-Einbettung}

Dieses Kapitel ist der kanonische Ort der Koordinaten- und CRS-Erklärung des
Teils Realität. In AIM definieren CRS und Koordinatensysteme den Einbettungs-
und Transformationskontext; sie definieren nicht die konstruktive
Alignment-Identität.

\section*{CRS sind keine Metadaten}

In der Alignment-Modellierung können CRS nicht als Hilfsmetadaten behandelt
werden. CRS-Annahmen beeinflussen metrische Interpretation,
Koordinatentransformationen, Projektionsverhalten und den Vergleich von Ziel-
und beobachteten Zuständen.

\begin{principle}[Strukturprinzip der CRS]
CRS sind ein struktureller Bestandteil des Realisierungs- und
Einbettungskontexts. Ihre Missachtung zerstört die Konsistenz der
ingenieurtechnischen Interpretation.
\end{principle}

Die CRS-Behandlung muss daher in Import, Laufzeitprojektion, Rekonstruktion
und realisierungsbewussten Vergleichsabläufen ausdrücklich erfolgen.


\section{Koordinatenreferenzsysteme}

\begin{definition}[Koordinatenreferenzsystem]
\index{Koordinatenreferenzsystem (CRS)}
Ein Koordinatenreferenzsystem (CRS) ist ein mathematischer Rahmen, der
räumliche Koordinaten auf Positionen in einer definierten World-Einbettung
abbildet.
\end{definition}

Ein CRS verbindet Annahmen zur Referenzfläche, Koordinatendefinitionen und
Transformationskonventionen.

\section{Einbettungskontext}

Eisenbahnprojekte verwenden häufig zugleich mehrere Systeme, darunter
ellipsoidbezogene Systeme, projizierte Systeme und lokale Ingenieurrahmen.
Diese Systeme stellen den Einbettungskontext für Geometrie und Messungen bereit.

\begin{principle}[Prinzip des CRS-Kontexts]
\index{Prinzip des CRS-Kontexts}
CRS gehören zum Realisierungs- und Einbettungskontext, nicht zur
konstruktiven Identität.
\end{principle}

\section{Transformationen und Konsistenz}

\begin{definition}[Koordinatentransformation]
\index{Transformation!Koordinaten}
Eine Koordinatentransformation ist eine Abbildung
\[
T:\R^n\rightarrow\R^n
\]
zwischen Koordinatensystemen.
\end{definition}

Die Behandlung von Transformationen muss ausdrücklich und prüfbar sein, da
kleine Transformationsinkonsistenzen große Interpretationsfehler in
krümmungs- oder projektionsbezogenen Größen hervorrufen können.

\section{Lokale Rahmen und numerische Stabilität}

\begin{definition}[Lokaler Rahmen]
\index{Bezugsrahmen!lokal}
Ein lokaler Rahmen ist ein Koordinatensystem, das nahe einem interessierenden
ingenieurtechnischen Gebiet verankert ist.
\end{definition}

Lokale Rahmen verbessern die numerische Stabilität von Laufzeitbewertung,
Projektion und Rekonstruktion und bewahren zugleich die Verbindung zum globalen
Referenzkontext.

\section{Stationierung gegenüber CRS-Koordinaten}

\index{Stationierung|see{Bezugssystem, Stationierung}}

CRS-Koordinaten drücken eine eingebettete Position aus. Stationierung drückt
eine Routenadresse entlang von Eisenbahnreferenzstrukturen aus. Dies sind
verschiedene Ebenen, die nicht vermengt werden dürfen.

\section{Bezug zur metrischen Realisierung}

\index{Realisierung!metrisch!Koordinatenbehandlung}

Koordinatenbehandlung verwendet metrische Realisierung
(\cref{chap:metric-realization}) und Realisierungsoperatoren; sie ersetzt sie
nicht. World-Einbettung und World--Track-Beziehungen bleiben
realisierungsabhängige Dienste.

\section{Verbindung zu AIM}

\begin{summary}
In AIM sind Koordinatensysteme und CRS ausdrücklicher Einbettungskontext für
realisierte und beobachtete Eisenbahnzustände.

Sie sind für Konsistenz und Interoperabilität wesentlich, bilden aber nicht
die Alignment-Identität.
\end{summary}
